\section{本章小结}
\label{sec:sqp-summary}

本章从圆周上的小问题出发，回答了章首提出的几个问题。SQP在当前点把目标近似到二阶、把约束近似到一阶，得到一个局部QP。目标的线性项来自$\nabla f$，曲率却要取Lagrange函数的Hessian，因为约束的弯曲只有通过乘子才折算得进模型。所以算法同时维护原始变量和乘子。

只含等式时，局部QP归结为一个KKT线性系统，可以整体分解，也可以分成法向段和零空间内的切向段。行满秩和约化Hessian正定保证解是模型的最小点。取精确Lagrange Hessian时，这个系统正是原问题KKT方程的Newton步。

含不等式时要决定哪些约束算数。有效集法维护工作集，步非零时用比值检验防止越界，步为零时用乘子的符号决定释放哪条约束。有效集法只是QP的一种解法，并不是SQP的一部分，QP也可以用内点法等方法求解。工作集属于当前的线性化QP，只有在约束资格、二阶充分条件和严格互补下，靠近解时才与原问题的活动集一致。

曲率方面，KKT矩阵可逆、约化Hessian正定、$\bm{B}$在全空间正定是三件不同的事。负曲率可以用移位、约化曲率修正或信赖域处理，它们分别改变模型或步长范围。拟Newton近似用同一组乘子下的Lagrange梯度差作割线，用阻尼保持正定，超线性收敛还需要沿步方向的切向曲率逐渐准确。

模型的步要放回原问题检验。罚函数把目标和违反量压成一个数，权重必须超过乘子，又不宜过大；过滤器并排比较两个指标，靠裕量、切换规则和恢复阶段才成为完整的机制。线性化约束无解时用弹性子问题或恢复阶段修复子问题，好步被弯曲约束的二阶违规挡住时用二阶修正修复候选步。恢复失败也不能证明原问题不可行。

最后要分清两类结论。解附近的快速收敛依赖约束资格、二阶充分条件、活动集的识别和全步长的接受；远离解时的可靠进展依赖接受规则和一致有界的曲率。两者条件不同，结论也不同，后者只给出KKT点，并不给出全局最优。全程都要分别检查QP的可行性与原非线性约束的可行性，前者只是后者在切平面上的影子。
